\cotefile{43}
\scanpage{17}
% archive sheet number: 16
% L1--3: top strip, crossed by long oblique strokes.
\diagram{\begin{tikzcd}[ampersand replacement=\&,column sep=large]
\deleted{C\times C'} \& \deleted{C\times C} \& \deleted{C}\arrow[r,no head] \& \deleted{\mathrm{Alb}} \& \deleted{\text{$0$ cycles}/G_m\text{-répartitions}}\\
\deleted{\cup} \& \deleted{\cup} \& \& \deleted{\mathrm{Pic}}\\
\deleted{D} \& \deleted{\Delta_C}
\end{tikzcd}}
% L4--8: upper left square and its labels.
\diagram{\begin{tikzcd}[ampersand replacement=\&,column sep=large,row sep=large]
\text{{\unsure{proalgébriques}}}\quad\prod_{\deleted{{\unsure{\mathrm{loc}}}}}G({\unsure{\hat{\mathcal O}_x|\hat O_x}})\arrow[d] \& 0\arrow[l]\arrow[d]\\
\substack{\text{{\unsure{pts. adéliques}}}\\=\text{{\unsure{répartitions}}}}\quad\prod_{\mathrm{loc}}G(\hat K_x) \& \underset{\substack{\text{{\unsure{fts rationnelles}}}\\=\text{rép. principales}}}{G(K_\xi)}\arrow[l]
\end{tikzcd}}
\margin{{\unsure{bicomplexe}} $DD(G)$}
\margin{Cohomologie $H^0(C,G)$ et $H^1(C,G)$}
% L9--13: second square; the aggregation signs have superposed vertical and zigzag strokes.
\diagram{\begin{tikzcd}[ampersand replacement=\&,column sep=large,row sep=large]
\text{``{\unsure{torsion}}''}\quad{\unsure{\sum|\prod|\coprod}}_{{\unsure{x|n}}}H_x^0(G)\arrow[d] \& 0\arrow[l]\arrow[d]\\
\text{diviseurs}\quad\sum_x H_x^1(G) \& \underset{\text{{\unsure{pts rationnels}} de $G$}}{H_\xi^0(G)}\arrow[l]
\end{tikzcd}}
\[
\begin{aligned}
H_{\mathrm I}(DD(G))&\\[-3pt]
={}&E_1^{**}(G)\Rightarrow H^*(C,G)
\end{aligned}
\]
en l'occurrence $H_x^0(G)=0$, d'où
\[
H^*(C,G)=\text{cohomologie}
\]
du {\unsure{complexe}}
\[
\underbrace{H_\xi^0(G)}_{\substack{G\text{-diviseurs}\\\text{principaux}}}
{\unsure{\longrightarrow|\text{---}}}
\underbrace{\prod_{{\unsure{x|\alpha}}} H_x^1(G)}_{G\text{-diviseurs}}
\]
% L14--19: lower left group, enclosed on the sheet by a curved outline.
\diagram{\begin{tikzpicture}\node[draw,rounded corners=7mm,inner sep=6pt]{\begin{tikzcd}[ampersand replacement=\&,column sep=large,row sep=large]
\underset{\phantom{a}}{\overset{\text{{\unsure{adèles entiers}}}}{\prod G(\hat{\mathcal O}_x)}}\arrow[d] \& 0\arrow[l]\arrow[d] \& G(C)=H^0({\unsure{C,C|G,C|C,G}}) \& 0\\
\underbrace{(\prod_{\mathrm{loc}}G(\hat K_x))/G(K_\xi)}_{\text{{\unsure{classes d'adèles}}}} \& 0\arrow[l] \& H^1(C,G) \& 0\\
H_{\mathrm{II}}DD(G) \& \& H_{\mathrm I}H_{\mathrm{II}}DD(G)
\end{tikzcd}};\end{tikzpicture}}
% L20--21.
qui peuvent aussi {\unsure{s'exprimer}} à l'aide de {\unsure{quotients locaux}} -----
\scanpage{18}
% archive sheet number: 17
% L1--3.
\diagram{\begin{tikzcd}[ampersand replacement=\&,column sep=large]
\prod J_x^{(0)}\arrow[r]\arrow[d,equal] \& J_\xi^*\\
J_1 \& J_0=J_{K/k}
\end{tikzcd}}
{\unsure{ou}}
\[
\deleted{{\unsure{\mathrm{gr}}}}\quad J_{C/k}^*\quad\text{jacobienne de }C
\]
% L4--12: crossed block. All its readable content is retained as deleted.
\deleted{La chose {\unsure{spéciale|précise}} et surprenante, c'est que l'on a aussi $J_{C/k}^*\simeq\underline{H}^1(C,G_m)$, {\unsure{et}} [si on {\unsure{prend}} le complexe $\operatorname{Hom}_{k\text{-}{\unsure{\mathrm{gr}}}}({\unsure{J_*|J^*}},G_m)$, on trouve}
\[
\deleted{K^*\longrightarrow\substack{\text{diviseurs}\\\text{{\unsure{ordinaires}}}}}
\]
\deleted{] et que $J_x^{(0)}(k)=G_m({\unsure{\hat{\mathcal O}_x|\hat O_x}})$.}

\deleted{Il faut aussi expliquer le rôle {\unsure{du|de}} $J_x^*$.}
% L13--19.
La chose {\unsure{non|un}} {\unsure{primordiale|principale}} : {\unsure{primo}}, c'est que le complexe
\[
J_1{\unsure{\longrightarrow|\text{---}}}J_0
\]
(ou plutôt $J_1(k){\unsure{\longrightarrow|\text{---}}}J_0(k)$) {\unsure{n'est}} autre que $H_{\mathrm{II}}DD(G_m)$, on peut donc, {\unsure{moyennant}} un abus de {\unsure{notations}},, écrire
\[
H_{\mathrm I}DD(G)=\operatorname{Hom}\bigl(\underline{H_{\mathrm{II}}DD}(G_m),G\bigr)
\]
\deleted{ce qui suggère que l'on ait}
\scanpage{19}
% archive sheet number: 18
% L1--3, upper diagrams.
\diagram{\begin{tikzcd}[ampersand replacement=\&,column sep=large,row sep=large]
\prod_x G({\unsure{\hat{\mathcal O}_x|\hat O_x}})\arrow[d] \& 0\arrow[l]\arrow[d] \& 0\arrow[d,rightsquigarrow] \& 0\arrow[l]\arrow[d,rightsquigarrow]\\
\prod_{\mathrm{local}}G(\hat K_x) \& G(K)\arrow[l] \& G\text{-{\unsure{diviseur|diviseurs}} sur }X \& G(K)\arrow[l]\\
\& \& H^1(X,G) \& H^0(X,G)
\end{tikzcd}}
% L4--9, small note below the first square.
i.e. on {\unsure{veut}} que {\unsure{ça corresponde|cela corresponde}} à un $G$-diviseur {\unsure{véritable}} sur $X$ [{\unsure{sinon donc|sinon dont}} le complexe {\unsure{continuerait-il}} ---]
% The horizontal rule between the two portions is a page separator.
\par\medskip\hrule\medskip
% L10--20: whole lower portion crossed by three oblique strokes.
\[\deleted{{\unsure{\sum|\prod|\coprod}}}\]
\[
\left\{\begin{aligned}
&\deleted{(f_x)_{x\in X}}\quad\deleted{\text{$G$-répartition}}\\
&\deleted{(\xi_D)_{D\in\mathrm{Div}\,\text{{\unsure{premiers}}}}}\qquad
\deleted{\xi_D\in J_{\underline{\mathcal O}_D/k}^*}
\end{aligned}\right.
\]
\[
\begin{aligned}
\deleted{x\in D}\qquad
\deleted{f_x\in G(K_x)}&\deleted{\longrightarrow\text{{\unsure{classe de}} }G(K_D)/\deleted{{\unsure{G(\mathcal O_D)}}}}\quad
\deleted{{\unsure{\simeq|\longrightarrow}}\ G(K_D)/G(\mathcal O_D)^{(1)}}\\
&\deleted{\hspace{4em}\simeq\operatorname{Hom}(J_{\underline{\mathcal O}_D/k}^*,G(k))}
\end{aligned}
\]
\deleted{Donc $(f_x,\xi_D)_x\in G(k)$ est défini.}

\deleted{Je dis que $\sum_{x\in|D|}(f_x,\xi_D)$ existe. En effet la question étant locale, on est {\unsure{ramené au cas}} où $(f_x)$ est formé d'unités,}
\scanpage{20}
% archive sheet number: 19
% L1--3: two sequences, with a crossed diagonal between them.
\diagram{\begin{tikzcd}[ampersand replacement=\&,column sep=large,row sep=large]
0\arrow[r] \& N\arrow[r]\arrow[dr,"?","{\times}" description] \& Z^0\arrow[r] \& \mathrm{Alb}^0\arrow[r] \&0\\
0\arrow[r,no head] \& \dfrac{\prod G({\unsure{\underline{\mathcal O}_x|\underline O_x}})}{G(C)}\arrow[r]\arrow[ur,no head] \& R_G^0/G(K)\arrow[r,no head] \& H^1(X,G)^0\arrow[r,no head] \&0
\end{tikzcd}}
% L4: overstruck superscript and discarded Alb.
\diagram{\begin{tikzcd}[ampersand replacement=\&,column sep=large]
X\arrow[r] \& \prod_D J_{\underline{\mathcal O}_D/k}^{\deleted{{\unsure{*}}}0}\arrow[r] \& J_{K/k}^0\arrow[r,no head,"\vert" description] \& \deleted{\mathrm{Alb}^0}J_{X/k}^0\arrow[r,no head] \&0
\end{tikzcd}}
% L5.
\diagram{\begin{tikzcd}[ampersand replacement=\&,column sep=large]
?\arrow[r] \& (\prod_D J_{\underline{\mathcal O}_D/k}^*)^0\arrow[r,"\vert" description] \& J_{X/k}^0\arrow[r,no head] \&0
\end{tikzcd}}
\par\hrule\medskip
% L6.
\diagram{\begin{tikzcd}[ampersand replacement=\&,column sep=large]
\text{\textcircled{?}}\arrow[r] \& \underset{\deleted{{\unsure{\mathrm{Diviseurs}}}}}{\prod J_{k(D)/k}^*}\arrow[r] \& J_{X/k}^*
\end{tikzcd}}
% L7--8.
\diagram{\begin{tikzcd}[ampersand replacement=\&,column sep=large,row sep=large]
\prod_{\mathrm{local}}J_{\underline{\mathcal O}_D/k}^*\arrow[r] \& J_{K/k}^*\\
\prod_{{\unsure{\mathrm{loc}}}}J_{\underline{\mathcal O}_D/k}^0\arrow[u]\arrow[r,no head] \& \prod J_{\underline{\mathcal O}_D/k}^0\arrow[u]
\end{tikzcd}}
\par\medskip\hrule\medskip
% L9--11: diamond; oblique quadrilateral around the left internal nodes.
\diagram{\begin{tikzcd}[ampersand replacement=\&,column sep=large,row sep=large,execute at end picture={\draw ([xshift=-2em]p20.west)--([yshift=2em]p20.north)--([xshift=2em]r20.east)--([xshift=1em,yshift=-1.5em]k20.south east)--([xshift=-2em]k20.south west)--([yshift=-.1em]c20.east)--cycle;}]
\&|[alias=p20]|\prod G(\underline{\mathcal O}_x)\arrow[dr,no head]\&\&J_G^0\arrow[dr,no head]\\
|[alias=c20]|G(C)\arrow[ur,no head]\arrow[dr]\&\&|[alias=r20]|R_G^0\arrow[ur,no head]\arrow[dr]\&\&H^1(X,G)^0\\
\&|[alias=k20]|G(K)\arrow[ur,no head]\&\&\mathrm{Div}_G^0\arrow[ur,no head]
\end{tikzcd}}
\par\medskip\hrule\medskip
% L12--14: final diamond. Question marks are manuscript text.
\diagram{\begin{tikzcd}[ampersand replacement=\&,column sep=large,row sep=large]
\&J_1=\prod J_{\underline{\mathcal O}_D/k}^0\arrow[dr]\&\&J_{K/k}^0\arrow[dr]\\
?\arrow[ur]\arrow[dr]\&\&(\prod_{\substack{D\\\mathrm{local}}}J_{\underline{\mathcal O}_D/k}^*)^0\arrow[ur]\arrow[dr]\&\&{\unsure{\mathscr J|J|\gamma}}_{X/k}^0\\
\&?\arrow[ur]\&\&({\unsure{\sum|\prod|\coprod}}_{D}\deleted{{\unsure{\mathrm{loc}}}}J_{k(D)/k}^*)^0\arrow[ur]
\end{tikzcd}}
\scanpage{21}
% archive sheet number: 20
% L1--3: upper left notes and upper right pairing.
\[
\begin{array}{cc}
A&A'\\
k&\deleted{{\unsure{K}}}L\\
\multicolumn{2}{c}{\overline{\deleted{{\unsure{G}}}\Gamma}}
\end{array}
\]
\diagram{\begin{tikzcd}[ampersand replacement=\&,column sep=large]
H^0(\Gamma,A_L)\times H^1(\Gamma,A'_L)\arrow[r] \& H^2(\Gamma,\deleted{{\unsure{K}}}L^*)\\
A_K\arrow[u,equal]
\end{tikzcd}}
\[
{\unsure{\operatorname{Ext}|\operatorname{Ext}^1}}_{{\unsure{\mathbb Z(\Gamma)/\mathbb Z|\mathbb Z[\Gamma]/\mathbb Z}}}(A_L,A'_L\deleted{{\unsure{*}}};L^*)
\]
% L4--8: left resolution sketch, crossed diagonals drawn without arrowheads.
\diagram{\begin{tikzcd}[ampersand replacement=\&,column sep=large]
A\&A'\\
A_1\arrow[u]\arrow[dr,no head]\&A'_1\arrow[u]\\
A_2\arrow[u]\arrow[ur,no head]\&A'_2\arrow[u]\\
0\arrow[u]\&0\arrow[u]
\end{tikzcd}}
\[
{\unsure{\mathcal L_*|\mathscr L_*}}\qquad \mathcal L'_*
\]
% L4--10: right product and composition diagram.
\diagram{\begin{tikzcd}[ampersand replacement=\&,row sep=large]
\begin{gathered}{\unsure{\operatorname{Ext}|\operatorname{Ext}^{1}|\operatorname{Ext}^{i}}}_\Gamma(\mathbb Z;A_L)\\\times\\{\unsure{\operatorname{Ext}|\operatorname{Ext}^{1}|\operatorname{Ext}^{j}}}_\Gamma(\mathbb Z,A'_L)\end{gathered}\arrow[d]\\
{\unsure{\operatorname{Ext}^{2}|\operatorname{Ext}^{1}|\operatorname{Ext}^{i+j}}}_\Gamma(\mathbb Z,\mathbb Z;A_L,A'_L)\times{\unsure{\operatorname{Ext}|\operatorname{Ext}^{1}}}_\Gamma(A_L,A'_L;L^*)\arrow[d]\\
\operatorname{Ext}_\Gamma^2(\mathbb Z,\mathbb Z;L^*)\arrow[d,"{\text{\uuline{{\unsure{cup}}}}}"]\\
{\unsure{\operatorname{Ex}|\operatorname{Ext}}}_\Gamma^2(\mathbb Z,L^*)
\end{tikzcd}}
% L9--13: left formulas.
\[
\operatorname{Hom}(\mathcal L_*,\mathcal L'_*;L^*)
\]
\[
\begin{aligned}
&H^1\bigl(\operatorname{Hom}(\mathcal L_*\otimes_{\mathbb Z}\mathcal L'_*,G)^\Gamma\bigr)\\
&H^1\bigl(\operatorname{Hom}({\unsure{M_*|\mathcal M_*}},\mathcal L_*)^\Gamma\bigr)\\
&H^1\bigl(\operatorname{Hom}(M_*,\mathcal L'_*)^\Gamma\bigr)
\end{aligned}
\]
\[
\begin{aligned}i+j+1&=2\\i+j&=1\end{aligned}
\qquad
\begin{array}{cc}\text{\textcircled{$i$}}&\text{\textcircled{$1-j$}}\\2&-1\end{array}
\]
% L14--17.
\diagram{\begin{tikzcd}[ampersand replacement=\&]
M_*\arrow[r,no head,"{\Gamma\text{-}{\unsure{\mathrm{hom}}}}"]\&\mathcal L_*\\
M_*\arrow[r,no head,"{\Gamma\text{-}\mathrm{hom}}"]\&\mathcal L'_*
\end{tikzcd}}
\[\deleted{P,Q}\]
\diagram{\begin{tikzcd}[ampersand replacement=\&,column sep=large]
\text{\uline{{\unsure{projectif}}}}\quad\mathcal L_*^P\otimes\mathcal L_*^Q\arrow[r,no head]\arrow[d] \& P\otimes Q\arrow[d,equal]\\
\text{\uline{résolution}}\quad\tilde{\mathcal L}_*^{P\otimes Q}\arrow[r]\&P\otimes Q
\end{tikzcd}}
% L18--20.
\diagram{\begin{tikzcd}[ampersand replacement=\&,column sep=large]
M_*\otimes_{\mathbb Z}M_*\arrow[r,"{\Gamma\text{-}\mathrm{hom}}"]\&\mathcal L_*\otimes_{\mathbb Z}\mathcal L'_*\arrow[r,"{\Gamma\text{-}\mathrm{hom}}"]\&G\\
M_*\arrow[u,"{\text{\uuline{{\unsure{cup}}}}\ 2}"]
\end{tikzcd}}
\[
\frac{1}{A_K/N_{L/k}A_L\ /\ \substack{\text{{\unsure{ss-groupe}}}\\\text{de }H^1(\Gamma,A'_L)}}
\]
% L21--22: bottom isolated items.
\[
K,L\qquad\mathbb Z\qquad A\quad P\quad\text{fibré principal sous }A
\]
\[
\deleted{H^1(}\qquad\underline{\operatorname{Pic}}_{P/S}
\]
\scanpage{22}
% archive sheet number: 21
% L1--3.
\diagram{\begin{tikzcd}[ampersand replacement=\&,column sep=large,row sep=large]
0\arrow[r,no head]\&G(K)/G(C)\arrow[r]\arrow[dr,no head]\&\mathrm{Div}_G^0\arrow[r,no head,"\vert" description]\&\deleted{{\unsure{D}}}H^1(X,G)^0\arrow[r,no head]\&0\\
\prod_{\substack{Z\text{ de}\\\text{codim }2}}J_{{\unsure{\underline{\mathcal O}_Z|\underline O_Z}}/k}\arrow[r]\&\underset{?}{\underline{\prod J_{\underline{\mathcal O}_D/k}^0}}\arrow[r,"u"']\arrow[ur,no head]\&J_{K/k}^0\arrow[r,"\vert" description]\&{\unsure{\mathscr J|J|\gamma}}_{X/k}^0\arrow[r,no head]\&0
\end{tikzcd}}
\margin{\uline{$X$ complet}}
% L4, cancelled line.
\[\deleted{\mathrm{Div}_G\qquad G(K)/G(C)}\]
% L5--6, two definitions and the small arrow from Im towards Ker.
\deleted{{\unsure{Div c Div}}}
\diagram{\begin{tikzcd}[ampersand replacement=\&,column sep=small]
(G\text{-}\mathrm{Div})=\operatorname{Hom}\bigl(\prod J_{\underline{\mathcal O}_D/k}^0/\&\operatorname{Im}(\prod_{Z\text{ cod., }2}J_{\underline{\mathcal O}_Z/k})\bigr)\arrow[d]\\
(G\text{-}\mathrm{Div})^0\underset{\text{{\unsure{pl.-isib.|pl.-isom}}}}{=}\operatorname{Hom}\bigl(\prod J_{\underline{\mathcal O}_D/k}^0\&\mid .\operatorname{Ker}u\bigr)
\end{tikzcd}}
% L7--9: boxed statements.
\[
\boxed{(G\text{-}\mathrm{Div})/(G\text{-}\mathrm{Div})^0\subset\operatorname{Hom}(H_1(J_*),G\deleted{{\unsure{(k)}}})}
\]
\[
\boxed{\text{Ex. Néron-Severi}\ \subset\operatorname{Hom}(H_1(J_*),G_m)}
\]
\fbox{A-t-on égalité ???}
\diagram{\begin{tikzcd}[ampersand replacement=\&]
0\arrow[r]\&A'\arrow[r,no head]\arrow[d,no head]\&A\arrow[r,no head]\arrow[d,no head]\&A''\arrow[r]\arrow[d,no head]\&0\\
\&\phantom{A'}\&\phantom A\&\phantom{A''}
\end{tikzcd}}
% L10--13: two lower sequences and crossing question-marked diagonals.
\diagram{\begin{tikzcd}[ampersand replacement=\&,column sep=large,row sep=large,execute at end picture={\draw ([xshift=.5em,yshift=1em]r22a.north east) .. controls ([xshift=4em]r22a.east) and ([xshift=4em]r22b.east) .. ([xshift=.5em,yshift=-1em]r22b.south east);}]
0\arrow[r,no head]\&\dfrac{\prod G(\underline{\mathcal O}_x)}{G(C)}\arrow[r]\arrow[dr,"?","\times" description]\&J_G\arrow[r]\&H^1(X,G)^0\arrow[r,no head]\&|[alias=r22a]|0\\
0\arrow[r,no head]\&?\arrow[r,no head]\arrow[ur,"?"]\&(\sum_D J_{k(D)/k}^*)^0\arrow[r,no head]\&\mathscr J_{X/k}^0\arrow[r,no head]\&|[alias=r22b]|0
\end{tikzcd}}
% slanted
\margin{Aucun {\unsure{accouplement}} {\unsure{convenable}} {\unsure{ne semble}} {\unsure{exister}} {\unsure{.....|....}}}
\[G(L)\]
{\unsure{Quel}} était l'accouplement {\unsure{-étique|-adique}} de Tate ???
\scanpage{23}
% archive sheet number: 22
% L1.
\[X\times Y\qquad A\quad{\unsure{\operatorname{Ext}|\operatorname{Ext}^1}}(A,G)\]
% L2--5.
$D$ \deleted{{\unsure{G}}} $G$-diviseur \inserted{\tikz[baseline=(w.base)]{\node[draw,rounded corners=12pt,inner sep=2pt] (w) {{\unsure{ordinaire|bilatéral}}};}} sur $X\times Y$, {\unsure{horizontal}} {\unsure{rel.}} à {\unsure{la projection}} $X\times Y\longrightarrow X$ transforme $0$-cycles de $X$ en $G$-diviseurs sur $Y$, et $0$-cycles du noyau d'Albanese en $G$-diviseurs principaux (?).

\deleted{{\unsure{Diviseur}}}
% L6--8.
\diagram{\begin{tikzcd}[ampersand replacement=\&]
X\times Y\arrow[dr]\\
\&{\unsure{\operatorname{Pic}^*(Y)}}
\end{tikzcd}}
\diagram{\begin{tikzcd}[ampersand replacement=\&]
J_{X/k}^*\arrow[r]\&\operatorname{Pic}_G^*(Y)
\end{tikzcd}}
% L9--14.
\diagram{\begin{tikzcd}[ampersand replacement=\&,column sep=large,row sep=large]
0\arrow[r,no head]\&\substack{\text{Noyau de}\\\text{Albanese}}\arrow[r,no head]\arrow[d]\&\underset{\substack{0\text{-cycles}\\\text{de deg }0}}{Z_X^0}\arrow[r]\arrow[d]\&\underset{\text{Albanese de }X}{J_{X/k}^0}\arrow[r,no head]\arrow[d]\&0\\
0\arrow[r,no head]\&G(K)/G(Y)\arrow[r]\&\underset{\substack{\text{{\unsure{au. alg.\textsuperscript{t}}}}\\\text{équiv. à }0}}{\mathrm{Div}^0(\deleted{{\unsure{X}}}Y,G)}\arrow[r]\&\underset{\substack{\text{classes de fibrés}\\\text{alg. équiv. à }0}}{\operatorname{Pic}^0(Y,G)}\arrow[r]\&0
\end{tikzcd}}
% L15, bottom sequence.
\diagram{\begin{tikzcd}[ampersand replacement=\&,column sep=large]
0\arrow[r,no head]\&\text{Noyau}\arrow[r]\&J_{K_Y/k}^0\arrow[r,no head]\&{\unsure{\mathscr J|J|\gamma}}_{Y/k}^0\arrow[r]\&0
\end{tikzcd}}
\scanpage{24}
% archive sheet number: 23
% L1, upper right title; the rest of the sheet is blank.
{\unsure{Théorie}} de Rosenlicht
\scanpage{25}
\mspage{1}
% archive sheet number: 24
% L1--2.
\uline{Théorie de Rosenlicht}

\textcircled{I} \uline{Existence des Picard pour les courbes}
% L3--9.
Soit $X$ un $Y$-schéma, avec

\begin{enumerate}[label=\alph*)]
\item $Y$ loc. noethérien
\item $X$ propre sur $Y$ ({\unsure{faut-il}} même : loc. projectif ?)
\item $X$ plat sur $Y$
\item Les fibres de $X$ sur $S$ sont des courbes algébriques \uline{absolument} \deleted{{\unsure{connexes}}} \deleted{{\unsure{irréductibles}}} \uline{réduites} et absolument \deleted{réduites}\inserted{connexes} \margin{[ce dernier pt n'est pas essentiel ---]}
\item Le nb de composantes \deleted{\illegible{1}} irréd. absolues des fibres de $X$ sur $Y$ est loc. constant sur $Y$
\end{enumerate}
% slanted
\margin{si on abandonne cette hypothèse, on {\unsure{trouve des|tombe aux}} {\unsure{phénomènes}} à la Néron ?!}
% L10.
Alors le {\unsure{$S$|$f$}}-schéma ${\unsure{\mathscr P_{X/Y}|\mathcal P_{X/Y}}}$ relatif existe, et est \uline{simple} sur $Y$.
% L11--17.
De plus, si le nb des composantes irréductibles absolues est $n$, il y a un revêtement \deleted{$Y$} étale $Y'$ de degré $n$ de $Y$ et si $\mathbb Z_{Y'/Y}$ est le schéma en groupes \inserted{$\prod_{Y'/Y}(\mathbb Z_{Y'})=\underline{\operatorname{Hom}}_Y(Y',\mathbb Z_Y)$} sur $Y$ défini par $Y'$ (qui sur $Y'$ devient isomorphe à $\mathbb Z^n$) on a une augmentation \deleted{{\unsure{surjective}} et {\unsure{propre}}} fidèlement plate
\margin{\uline{{\unsure{références}}}}
\[
\mathscr P_{X/Y}\longrightarrow\mathbb Z_{Y'/Y}\quad\deleted{{\unsure{\infty}}}
\]
dont le noyau sera noté $\mathscr P_{X/Y}^c$, d'où la suite exacte
\[
0\longrightarrow\mathscr P_{X/Y}^c\longrightarrow\mathscr P_{X/Y}\longrightarrow\mathbb Z_{Y'/Y}{\unsure{\longrightarrow|\text{---}}}0
\]
% L18--22.
{\unsure{Noter}} que l'on a un \uline{homomorphisme {\unsure{canon.|canonique}}}
\[
\mathbb Z_{Y'/Y}\longrightarrow\mathbb Z_Y
\]
donc une augmentation
\[
\mathscr P_{X/Y}\longrightarrow\mathbb Z_Y
\]
\scanpage{26}
% archive sheet number: 25
% L1--3
Les fibres de ${\unsure{\mathscr P|\mathcal P}}^0_{X/S}$ sont connexes, et propres si celles de $X$ sont \uline{simples}. Soit \deleted{{\unsure{schéma}}} $X-S$ l’{\unsure{ens.|ouv.}} des pts de $X$ qui sont {\unsure{non}} simples sur $Y$. On va définir un morphisme
% L4
\[
\boxed{X-S\longrightarrow\mathscr P^1_{X/S}}.
\]
% L5--8
où l’exposant $1$ indique qu’on est dans la composante d’augmentation $1$. {\unsure{Je reviens|Il revient}} {\unsure{à en donner}} \uline{fonctoriellement} \inserted{pour le changement de base} une application
\deleted{{\unsure{application}}}
\[
{\unsure{\Gamma|\mathscr T}}(X-S/Y)\longrightarrow\mathscr P^1(X/S)
\]
% L9--13
or on note qu’une section de $X-S$ sur $Y$ \uline{définit un diviseur sur} $X-S$ {\unsure{(transverse à $X/Y$)}} d’où un fibré inversible, dont on constate que le \deleted{{\unsure{diviseur}}} \inserted{degré} est $+1$ [car on peut supposer ${\unsure{Y|S}}=\operatorname{Spec}k$, {\unsure{le degré clair}}].
% margin note written slanted, rising toward the text
\margin{{\unsure{vérifier|référence}}??}

% L14--23
II) Partons d’un $X$ \uline{propre} et \uline{simple} sur ${\unsure{Y|S}}$, dont les fibres sont de dim $1$, \deleted{{\unsure{Soit}} $X_i$} {\unsure{(et abs! connexes)}} (donc abs! irréductibles). Soit $S$ un sous-schéma \inserted{fermé} de \deleted{$X$ plat sur $Y$ et} ne contenant aucune fibre de $X/{\unsure{Y|S}}$, donc \uline{fini} sur ${\unsure{Y|S}}$, et plat sur $Y$. Soit ${\unsure{\Xi|\mathscr I|\mathcal I}}$ le faisceau d’idéaux définissant $Y$. Alors on prouve facilement que $\Xi$ est \uline{localement libre} [c’est purement local : un faisceau coh. sur $X$ plat rel. à $Y$ et qui {\unsure{est inversible|est un module}} dans fibres \inserted{{\unsure{est}}} est plat sur $X$.---]
% margin note written slanted, read after rotating clockwise
\margin{{\unsure{fibres abs! connexes est une condition}} {\unsure{peut-être}} \uline{{\unsure{vide}}}.}
\scanpage{27}
% archive sheet number: 26
% L1--4
donc \deleted{{\unsure{il s’ensuit}}} \inserted{{\unsure{il s’ensuit}}} \deleted{{\unsure{fac.}}} {\unsure{facile|facilement}} que les ${\unsure{\Xi|\mathscr I|\mathcal I}}^{{\unsure{-n|n|-m|m}}}$ sont tels que ${\unsure{\underline{\mathcal O}|\underline O}}_X/\Xi^{-n}$ sont encore plats sur $Y$. Soit $\deleted{{\unsure{Y}}}S_n$ le sous-schéma de $X$ défini par $\Xi^{-n}$, et soit $X_n=X\amalg_{S_n}Y$ le $\deleted{S}Y$-schéma obtenu en ``{\unsure{contractant}} $S_n$ en $S$''.
% L5--6
\deleted{{\unsure{Soit}}} \deleted{{\unsure{de sorte}} $X_{{\unsure{S|\xi}}}=X-S$.} Donc $X_0=X$, et les $X_i$ forment un \uline{système} \uline{\deleted{projectif}\inserted{inductif}}.
% left, cancelled square and replacement square
\diagram{\begin{tikzcd}[ampersand replacement=\&]
\deleted{{\unsure{S}}_n}\arrow[r,no head]\arrow[d]\arrow[dr,no head] \& \deleted X\arrow[d] \& S_n\arrow[r,"i_n"]\arrow[d,"f i_n"'] \& X\arrow[d,"p_n"]\\
\deleted Y\arrow[r,no head] \& \deleted{X_n} \& Y\arrow[r,"\varepsilon_n"'] \& X_n
\end{tikzcd}}
\diagram{\begin{tikzcd}[ampersand replacement=\&]
\cdots\& X_2\arrow[l,no head]\&X_1\arrow[l]\&X_0=X\arrow[l]
\end{tikzcd}}
% L7--14
Donc les ${\unsure{\mathscr P|\mathcal P}}_{X_n/S}$ forment un système \deleted{\uline{inductif}} projectif, d’où un pro-($S$-préschéma en groupes) noté $\mathscr P_{S,X/Y}$, et qui ne dépend que de la \uline{partie fermée} $S$ de $X/Y$, comme on constate facilement (du moins, si $Y$ est noethérien) --- cel{\unsure{a}} {\unsure{résultera|résulte}} d’ailleurs {\unsure{de}} la ``construction'' qu’on en {\unsure{donnera|donne}} plus bas.

% L15--17
Les morphismes canoniques $X-S\longrightarrow\mathscr P^1_{X_n/Y}$ définissent à la limite un morphisme
\[
\boxed{{\unsure{\varphi|\phi}}_S:X-S\longrightarrow\mathscr P^1_{S,X/Y}{\unsure{\subset|\in}}\mathscr P_{S,X/Y}}.
\]
\scanpage{28}
% archive sheet number: 27
% L1--7; vertical emphasis stroke beside the statement
\noindent\vrule\hspace{.6em}\begin{minipage}{.93\linewidth}
\uline{Théorème} Soit $G$ un $Y$-schéma en groupes \inserted{{\unsure{commu.}}} de type fini \deleted{sur ${\unsure{Y|S}}$}, et $P^1$ un fibré principal homogène sous $G$ [devenant trivial par extension fid. plate quasi-compacte{\unsure{---}}], d’où un schéma en groupes $\mathbb Z$-augmenté $P$, avec $P^{(1)}=P^1$, $P^{(0)}=G$. Considérons l’application canonique $u{\unsure{\longrightarrow|\mapsto}}u\circ{\unsure{\varphi|\phi}}_S$
\end{minipage}\par
\[
\operatorname{Hom}_{\substack{Y\text{-schémas en groupes}\\\mathbb Z\text{-augmentés}}}
({\unsure{\mathscr P|\mathcal P}}_{S,X/Y},P)\longrightarrow\operatorname{Hom}_Y(X-S,P^1)
\]
% L8--9
\uline{Cette application est bijective}.

(Th. de Rosenlicht)

% L10
\uline{Corollaire 1}
\[
\operatorname{Hom}_{Y\text{-schémas en groupes}}(\mathscr P_{S,X/Y},G)\xrightarrow{\sim}\operatorname{Hom}_Y(X-S,G)
\]
% L11--15; vertical emphasis stroke beside lower paragraph
\uline{Corollaire 2}
\noindent\vrule\hspace{.6em}\begin{minipage}{.93\linewidth}
\inserted{Considérons} Le \deleted{{\unsure{homomorphisme de}}} \inserted{foncteur} $\varphi_S^*$, de la catégorie des schémas \uline{extensions} \inserted{(en groupes commutatifs)} de $P_{S,X/Y}$ par des $Y$-schémas en groupes commutatifs de type fini $G$ sur $Y$, dans la catégorie \deleted{{\unsure{qui}}}
\end{minipage}\par
% margin note written slanted, read after rotating clockwise
\margin{il faut faire {\unsure{attention aux définitions}} de ``fibré'' et ``extension'' ? $G$ plat sur ${\unsure{Y|S}}$??}
\scanpage{29}
% archive sheet number: 28
% L1--3, continuation of Corollaire 2
\noindent\vrule\hspace{.6em}\begin{minipage}{.93\linewidth}
des fibrés principaux homogènes sur $X-S$ de groupes de tels $G$. Ce foncteur est pleinement fidèle.
\end{minipage}\par

% L4--9
Soient $Q,Q'$ deux telles extensions, de groupes $G,G'$. Un hom. de $Q'$ en ${\unsure{Q'|Q}}$ est \inserted{un couple formé d’un hom. $u:G'\longrightarrow G$ et d’} un{\unsure{e}} \deleted{isom. de $Q''=Q\times_{G'}G$ avec} \deleted{{\unsure{élément de}} $\operatorname{Hom}_{{\unsure{S,u}}}({\unsure{Q'|G'}},Q)$, {\unsure{avec}}} \deleted{qui est un fibré} \inserted{{\unsure{isom. de}} $Q''={\unsure{Q'\times_{G'}G|Q\times_{G'}G}}$} $Q'$ \deleted{{\unsure{ou}}} (on {\unsure{entend}} par extension de groupe \uline{donné} $G'$). Il y en a {\unsure{ssi}} ${\unsure{Q'|Q}}$ et $Q''$ sont isomorphes. Or
% margin note written slanted, rising to the right
\margin{il faut {\unsure{pouvoir}} définir les fibrés associés, Attention!}
% left-middle diagram
\diagram{\begin{tikzcd}[ampersand replacement=\&]
\&Q\arrow[d,no head]\&Q''\\
X-S\arrow[ur,dashed]\arrow[r]\&{\unsure{\mathscr P|\mathcal P}}_{S,X/S}
\end{tikzcd}}
% L10 displayed, overwritten letters retained
\[
{\unsure{\operatorname{Ext}|\operatorname{Ext}^1}}({\unsure{\mathscr P|\mathcal P}}_{{\unsure{S}}},G)\longrightarrow\deleted{{\unsure{\operatorname{Ext}}}}{\unsure{H|\mathrm H}}^1(X-S,G)
\]
% L11--17
étant un hom. de groupes, pour prouver qu’il est injectif (i.e. $Q$ et $Q'$ isom. ssi ${\unsure{\varphi|\phi}}^*(Q)$ et $\varphi^*(Q'')$ le sont) on est ramené à prouver que le noyau est nul, i.e. que si $\varphi^*(Q)$ est trivial, alors $Q$ l’est. Cela résulte \deleted{facilement} \inserted{{\unsure{aussitôt}}} du th. de Rosenlicht.
% margin note written slanted, read in its own direction
\margin{Attention, il faut bien préciser les définitions {\unsure{d’opérations d’addition}} de fibrés etc}
% L18, unfinished sentence at foot
Comme \inserted{$\varphi^*$ est fidèle car} ({\unsure{le groupe étant donné}}), \deleted{{\unsure{est}}} {\unsure{on|en}}
\scanpage{30}
% archive sheet number: 29
% L1--3, continuation
est ainsi ramené à prouver que
\[
\operatorname{Hom}(Q,Q)\longrightarrow\operatorname{Hom}({\unsure{\varphi|\phi}}^*(Q),\varphi^*(Q))
\]
est bijectif. Or les deux s’identifient {\unsure{resp!|resp.}} à
% L4--5
\[
\operatorname{Hom}_{\substack{\text{grps}\\\text{sur }{\unsure{Y|S}}}}({\unsure{\mathscr P|\mathcal P}}_S,G)
\quad\text{et}\quad
\deleted{{\unsure{E|H}}}\operatorname{Hom}_Y(X-S,G)
\]
donc sont isomorphes.

% L6--9
En particulier

\noindent\vrule\hspace{.6em}\begin{minipage}{.93\linewidth}
\uline{Corollaire 3} Si $G$ est un schéma en groupes commu. de type fini sur $Y$, on a un hom. canonique \uuline{injectif}
\end{minipage}\par
\[
\operatorname{Ext}_{\substack{\text{schémas en groupes}\\\text{sur }{\unsure{Y|S}}}}(\mathscr P_{S,X/S},G)\longrightarrow H^1(X-S,G)
\]
% margin note written slanted
\margin{{\unsure{préciser la signification}}{\unsure{...}}}
% L10--15
Il \uline{semble} que c’est \uline{{\unsure{bijectif|surjectif}}} chaque fois que $S\longrightarrow Y$ est \uline{surjectif}, i.e. que les fibres de $X/Y$ sont non propres, donc \uline{affines}. Dans le cas où $S=\varnothing$, $G=G_m$, on ne {\unsure{sera|serait}} déjà pas surjectif.
\scanpage{31}
% archive sheet number: 30
% L1--7
Il semble néanmoins que dans le cas général, pour prouver qu’un fibré principal $E$ sur $X-S$ provient d’une \deleted{fibré principal} extension de ${\unsure{\mathscr P|\mathcal P}}_S$ par $G$, il suffit (utilisant {\unsure{la théorie|le théorème}} de la descente) de le vérifier en {\unsure{remplaçant}} $S$ par les spectres des \uline{complétés} de ses anneaux locaux.
% L8--10
De plus, utilisant la géométrie formelle [lorsque $G$ est bien plongé dans le groupe \deleted{{\unsure{H|\mathrm H}}} linéaire] on est ramené à le vérifier \uline{formellement}.
% L11--14
Cela revient au cas d’un {\unsure{anneau}} de base {\unsure{Artinien}}. Travaillant de proche en proche, on est ramené à le vérifier {\unsure{pour les courbes}} \uline{{\unsure{sur|pour}}} les corps résiduels de $S$.
% L15--19
Utilisant à nouveau la descente on serait {\unsure{amené|ramené}} à le vérifier à nouveau sur les clôtures algébriques des corps résiduels{\unsure{...---}} On peut sans doute s’en tirer plus {\unsure{astucieusement}}.
